\begin{solapartat}
\punts{0,25} Tal com explica l'enunciat, un ull miop crea la imatge d'objectes llunyans per davant de la retina. Per tant necessitem posar davant de l'ull una lent divergent que faci divergir (separar, obrir) els feixos de llum, i així l'ull formi la imatge a la retina i no per davant d'aquesta.

\punts{0,25} Com que l'ull sols pot veure objectes situats a una distància màxima de 4 metres, la lent divergent haurà de crear la imatge d'un objecte llunyà (a l'infinit $s = -\infty$, o $s'$ molt gran) a 4~metres per davant de l'ull $s' = -4\ \mathrm{m}$.

\punts{0,25} Per trobar la focal utilitzem altra vegada l'equació de la lent: $\dfrac{1}{s'} - \dfrac{1}{s} = \dfrac{1}{f'}$

\punts{0,25} Introduïm els valors: $\dfrac{1}{-4} - \dfrac{1}{-\infty} = \dfrac{1}{f'}$, per tant la focal de la lent és $f' = -4\ \mathrm{m}$.

\punts{0,25} La potència d'una lent (en diòptries) és la inversa de la focal (en metres). Per tant:
\[ P = \frac{1}{f'} = \frac{1}{-4} = -0,25\ D. \]
\end{solapartat}

\begin{solapartat}
\punts{0,25} Per trobar la distància a la que es forma la imatge creada per la lent correctora de l'arbre que es troba a 12 metres, $s = -12\ \mathrm{m}$ utilitzem l'equació de la lent: $\dfrac{1}{s'} - \dfrac{1}{s} = \dfrac{1}{f'}$ amb el valor de la focal que ja hem trobat: $\dfrac{1}{s'} - \dfrac{1}{-12} = \dfrac{1}{-4}$

\punts{0,25} Per tant $\dfrac{1}{s'} = \dfrac{1}{-4} + \dfrac{1}{-12} = \dfrac{-3-1}{12} = \dfrac{-4}{12} = -\dfrac{1}{3}$ i per tant $s' = -3\ \mathrm{m}$.

\punts{0,75} El raig que surt de l'arbre i és paral·lel a l'eix passa pel punt focal imatge $f'$, situat a $-4\ \mathrm{m}$ de la lent. El raig que passa pel centre de la lent no es desvia. El raig que surt de l'arbre i passa per $f$, situat a 4~m surt paral·lel a l'eix. Sols són necessaris dos rajos per determinar la posició de la imatge.

\figura[0.95]{sol_fig1.png}
\end{solapartat}
